%─────────────────────────────────────────────────────────────────────────────
% CHAPTER 2 — State of the Art
% Talan Tunisia — PFE Report
%─────────────────────────────────────────────────────────────────────────────

\chapter{State of the Art}
\thispagestyle{fancy}

%─────────────────────────────────────────────────────────────────────────────
\section{Introduction}
%─────────────────────────────────────────────────────────────────────────────

The Intelligent Enterprise Strategic Assistant sits at the intersection of
several rapidly evolving research fields: large language models, multi-agent
AI architectures, retrieval-augmented generation, knowledge representation,
graph-based reasoning, graph neural networks, and enterprise information
systems integration. This chapter reviews the state of the art in each of
these areas, identifies the key contributions that directly inform the design
of the proposed system, and positions the project relative to existing
academic and industrial work.

%─────────────────────────────────────────────────────────────────────────────
\section{Enterprise Information Systems and the Integration Challenge}
%─────────────────────────────────────────────────────────────────────────────

\subsection{ERP, CRM, and HR Systems}

Modern enterprises rely on three classes of domain-specialized information
systems. Enterprise Resource Planning (ERP) systems manage financial flows,
procurement, and supply chain operations. Customer Relationship Management
(CRM) platforms centralize sales pipelines, client interactions, and
commercial opportunities. Human Resource Information Systems (HRIS) handle
employee records, contracts, skills, leave management, and project
assignments. While each system excels in its own domain, they are designed
as independent silos with distinct data models, APIs, and access
paradigms~\cite{laudon2020mis}.

\subsection{The Fragmentation Problem}

The coexistence of these isolated platforms creates a structural integration
challenge: cross-domain queries — such as identifying which employees are
assigned to projects for clients showing declining revenue — require manual
correlation across systems that share no common semantic layer.
Enterprise Application Integration (EAI) and more recent data mesh
architectures~\cite{dehghani2022datamesh} attempt to address data movement
at the infrastructure level, but they do not provide an intelligent
interaction surface. Decision-makers are thus left to reconstruct a unified
view of the organization manually, a process that is slow, error-prone, and
does not scale with organizational complexity.

\subsection{Text-to-SQL for Natural Language Database Access}

Text-to-SQL research addresses the challenge of translating natural language
queries into executable SQL statements, enabling non-technical users to
interrogate relational databases directly. Recent LLM-based approaches have
demonstrated competitive performance on standard benchmarks such as
Spider~\cite{yu2018spider} and BIRD~\cite{li2023bird}, showing that
instruction-tuned models can reliably generate correct SQL across diverse
schemas when provided with schema context in the prompt. This capability is
a prerequisite for any conversational assistant that aims to provide
domain-neutral access to enterprise relational databases.

\subsection{Role-Based Analytics and Dashboards}

Business intelligence research has long recognized that different
organizational roles require fundamentally different views of the same
underlying data~\cite{watson2017bi}. An employee, a project manager, and a
C-suite executive have distinct information needs, decision horizons, and
risk tolerances. Role-based dashboards — adaptive interfaces that surface
the most relevant metrics and alerts for a given user profile — have become
a standard pattern in enterprise analytics. The present system extends this
concept by combining role-aware visual dashboards with a conversational
agent layer, allowing users to seamlessly transition between structured
visualization and natural-language exploration within the same platform.

%─────────────────────────────────────────────────────────────────────────────
\section{Large Language Models}
%─────────────────────────────────────────────────────────────────────────────

\subsection{From Transformers to Foundation Models}

The Transformer architecture, introduced by Vaswani et al.~\cite{vaswani2017attention},
established the technical foundation for all modern large language models (LLMs).
By replacing recurrence with self-attention mechanisms, Transformers enabled
the parallel processing of sequences and the capture of long-range dependencies
at scale. Subsequent scaling laws demonstrated that model capability improves
predictably with parameter count and training data volume~\cite{kaplan2020scaling}.

The GPT series~\cite{brown2020gpt3} demonstrated that sufficiently large models,
trained on diverse text corpora, acquire broad task-solving abilities without
explicit fine-tuning — a property termed \textit{in-context learning}. This
breakthrough established the \textbf{foundation model} paradigm, in which a
single pre-trained model serves as the basis for a wide range of downstream
applications through prompting, instruction tuning, or parameter-efficient
adaptation~\cite{bommasani2021foundation}.

\subsection{Reasoning and Tool Use}

Despite their generative power, early LLMs suffered from factual
hallucinations and limited step-by-step reasoning. Chain-of-thought
prompting~\cite{wei2022chainofthought} addressed the latter by instructing
models to verbalize intermediate reasoning steps before producing a final
answer, yielding substantial gains on complex arithmetic and commonsense
tasks.

The ReAct framework~\cite{yao2023react} further extended this capability by
interleaving reasoning traces with tool-calling actions, enabling LLMs to
interact with external environments — search engines, calculators, databases —
in a controlled loop. This \textit{reason-then-act} paradigm is the direct
conceptual precursor to modern AI agent architectures and underlies the
design of the domain agents in the present system.

\subsection{LLM-Based Task Automation}

Beyond question answering, LLMs have demonstrated the ability to automate
structured operational tasks by translating natural language intent into
concrete actions: filling forms, drafting documents, generating reports, and
managing calendar or workflow entries~\cite{shen2023hugginggpt}. This
capability is particularly relevant in enterprise settings, where repetitive
tasks such as leave request processing, email drafting, and meeting
preparation consume significant managerial time. By embedding task automation
directly into the agent layer, AI assistants can serve as active operational
collaborators rather than passive information retrieval systems.

%─────────────────────────────────────────────────────────────────────────────
\section{Multi-Agent AI Architectures}
%─────────────────────────────────────────────────────────────────────────────

\subsection{Foundations of AI Agents}

An AI agent is an autonomous software component that perceives its
environment through context (user messages, tool results, conversation
history), reasons over it via an LLM, and executes actions through a set of
registered tools. The \textit{Reflexion} framework~\cite{shinn2023reflexion}
demonstrated that agents capable of verbal self-reflection on past failures
achieve markedly higher task completion rates, introducing the notion of
\textit{episodic memory} at the agent level.

\subsection{Multi-Agent Orchestration}

Complex enterprise tasks — those requiring simultaneous access to HR, CRM,
ERP, and market data — exceed the reliable scope of a single agent.
Multi-agent systems address this by decomposing complexity: a central
orchestrator interprets user intent and delegates subtasks to specialized
agents, each equipped with domain-specific tools and knowledge.
Park et al.~\cite{park2023generative} showed that collections of LLM-backed
agents can sustain coherent behavior and emergent collaboration dynamics,
providing theoretical grounding for agent cooperation in organizational
simulations.

AutoGen~\cite{wu2023autogen} demonstrated that conversational multi-agent
frameworks, in which agents communicate by exchanging structured messages,
can solve complex multi-step tasks across domains without hardcoded
workflows. This contrasts with earlier approaches that relied on
pre-specified agent graphs with fixed topologies and no dynamic routing.

\subsection{LangGraph for Stateful Agent Orchestration}

LangGraph~\cite{langgraph2024} extends the LangChain ecosystem with a
graph-based execution model in which agents and tools are nodes, and
conditional transitions define the control flow. Unlike linear chains,
LangGraph natively supports cycles, branching, and shared state
persistence — properties essential for multi-turn conversations that span
multiple domains. Its explicit state machine semantics allow for
deterministic routing of user intent to the appropriate domain agent, as
well as cross-domain orchestration when queries require evidence from more
than one enterprise system. In the present project, LangGraph serves as the
central orchestration engine binding the HR, CRM, ERP, Market Analysis,
and Competitive Intelligence agents into a coherent, stateful assistive
system.

%─────────────────────────────────────────────────────────────────────────────
\section{Retrieval-Augmented Generation}
%─────────────────────────────────────────────────────────────────────────────

\subsection{Motivation and Principle}

LLMs encode world knowledge in their parameters during training, but this
knowledge is static: it cannot reflect organizational documents, internal
policies, or events that post-date the training cut-off. Retrieval-Augmented
Generation (RAG), introduced by Lewis et al.~\cite{lewis2020rag}, resolves
this limitation by dynamically retrieving relevant passages from an external
corpus at inference time and injecting them into the model's context window.
The model can then ground its response in retrieved evidence rather than
relying solely on parametric memory, substantially reducing hallucinations
on knowledge-intensive tasks.

\subsection{Retrieval Pipeline}

A standard RAG pipeline comprises three stages~\cite{gao2023rag}:

\begin{enumerate}
  \item \textbf{Indexing.} Source documents are segmented into chunks,
        encoded into dense vector representations by an embedding model,
        and stored in a vector database alongside metadata.
  \item \textbf{Retrieval.} At query time, the user question is embedded
        using the same model, and the $k$ most semantically similar chunks
        are retrieved by approximate nearest-neighbor search.
  \item \textbf{Generation.} The retrieved passages are prepended to the
        prompt, providing the LLM with factual context from which to
        synthesize a grounded, cited response.
\end{enumerate}

\textbf{Hybrid search} — combining dense cosine similarity with sparse BM25
lexical scoring~\cite{robertson2009bm25} — has been shown to outperform
either method alone, particularly when queries contain specific technical
terms or identifiers that semantic embeddings may dilute. In the present
system, hybrid search is applied over the corpus of internal documents
(company policies, procedures, and reports) to maximize retrieval recall
for employee and manager queries about organizational rules and guidelines.

\subsection{Advanced RAG Patterns}

Beyond the vanilla pipeline, several extensions have proven effective in
enterprise settings~\cite{gao2023rag}. \textit{Contextual compression}
filters retrieved passages to retain only sentences relevant to the query,
reducing context noise. \textit{Reranking} applies a cross-encoder model
to reorder the initial retrieval results by relevance before generation.
\textit{Self-RAG}~\cite{asai2023selfrag} introduces a reflective mechanism
in which the model decides when retrieval is needed and critiques its own
retrieved evidence, improving both efficiency and faithfulness. These
patterns collectively inform the RAG module design of the present system,
where retrieval is triggered selectively based on the nature of the user's
query.

%─────────────────────────────────────────────────────────────────────────────
\section{Knowledge Graphs}
%─────────────────────────────────────────────────────────────────────────────

\subsection{Definition and Structure}

A knowledge graph (KG) is a directed, labeled graph in which nodes represent
entities (persons, organizations, products, events) and edges represent typed
relationships between them~\cite{hogan2021knowledge}. Formally, a KG is a
set of subject–predicate–object triples $(s, p, o)$, optionally annotated
with provenance metadata. This structure enables multi-hop reasoning:
inferring facts that are not explicitly stored by traversing chains of
relationships across the graph. Enterprise KGs model the company's
operational reality — who works on what project, which client is in which
sector, which financial indicator correlates with which business outcome —
in a form that is both machine-queryable and human-interpretable~\cite{noy2019industry}.

\subsection{Graph Databases: Neo4j}

Neo4j~\cite{neo4j2024} is the leading native graph database, storing data
as nodes and typed relationships with properties, and exposing them through
the Cypher query language. Its architecture is optimized for
traversal-heavy queries — such as finding the shortest impact path between
a market event and a set of affected projects — that would require costly
multi-join operations in a relational model. In the present system, Neo4j
hosts the enterprise knowledge graph, connecting HR, CRM, and ERP entities
with external market events to enable cross-domain reasoning, world model
visualization, and impact propagation analysis.

\subsection{Temporal Knowledge Graphs}

Standard knowledge graphs represent a static snapshot of the world. Temporal
knowledge graphs (TKGs) extend this model by associating facts with time
intervals or timestamps: $(s, p, o, t)$, allowing the graph to encode the
evolution of relationships and events over time~\cite{trivedi2019dyrep}.
TKGs are particularly relevant for market intelligence applications, where
the chronological ordering of events — a competitor's product launch,
a regulatory announcement, a market downturn — determines their causal
interpretation and strategic significance. In the present system, the
knowledge graph is enriched with temporally annotated market events, enabling
the GNN propagation layer to reason about the sequential impact of external
signals on the enterprise's internal network.

\subsection{Interactive Graph Visualization and World Models}

Graph visualization tools such as D3.js and Cytoscape.js enable interactive
exploration of complex entity-relationship networks directly in the browser.
The concept of an enterprise \textit{world model} — a unified,
graph-based representation of all the organization's actors, processes,
assets, and their interrelationships — has been proposed as a foundation
for organizational situational awareness~\cite{noy2019industry}. By
rendering this graph interactively, users can navigate between internal
entities (employees, projects, clients) and external market entities
(competitors, sectors, events) within a single visual canvas, revealing
structural dependencies that tabular reports cannot expose.

\subsection{GraphRAG}

GraphRAG~\cite{edge2024graphrag}, proposed by Microsoft Research, extends
standard RAG by grounding LLM responses in a structured knowledge graph
rather than raw document chunks. Entity-centric retrieval traverses graph
neighborhoods to surface relational context — co-workers, past projects,
client history — that flat vector search cannot capture. This approach is
particularly valuable in enterprise settings where the relationships between
entities, not just isolated documents, carry the most strategic information.

%─────────────────────────────────────────────────────────────────────────────
\section{Graph Neural Networks for Enterprise Intelligence}
%─────────────────────────────────────────────────────────────────────────────

\subsection{Graph Convolutional Networks}

Graph Neural Networks (GNNs) generalize deep learning to graph-structured
data by learning node representations through iterative message passing:
each node aggregates feature vectors from its neighbors, enabling the model
to capture structural patterns at increasing scales~\cite{kipf2017gcn}.
The Graph Convolutional Network (GCN) by Kipf and Welling~\cite{kipf2017gcn}
formalized this spectral approach and achieved strong performance on
semi-supervised node classification tasks.

\subsection{Graph Attention Networks}

Graph Attention Networks (GAT)~\cite{velickovic2018gat} introduced
attention-weighted aggregation, allowing each node to assign learned
importance scores to its neighbors rather than treating all edges equally.
This is particularly relevant in heterogeneous enterprise graphs where
some relationships — a client responsible for a significant share of revenue,
a key employee bridging multiple strategic projects — are far more
consequential than others and should contribute more heavily to impact
estimates.

\subsection{Inductive Learning with GraphSAGE}

GraphSAGE~\cite{hamilton2017graphsage} addressed the scalability limitation
of transductive GNNs by learning aggregation functions that generalize to
unseen nodes — a critical property in a dynamic enterprise graph where new
employees, clients, and market events are continuously added. In the present
project, GNN-based propagation is applied to the temporal knowledge graph
to assess how external market signals (competitor announcements, sectoral
downturns, regulatory changes) cascade through the organization's network
of clients, projects, and human resources, producing quantified risk scores
and opportunity indicators for strategic decision-makers.

%─────────────────────────────────────────────────────────────────────────────
\section{Market Intelligence and Competitive Analysis}
%─────────────────────────────────────────────────────────────────────────────

\subsection{Competitive Intelligence as a Discipline}

Competitive intelligence (CI) is the systematic process of collecting,
analyzing, and distributing actionable information about the competitive
environment to support strategic decision-making~\cite{gilad1989ci,
porter1980competitive}. The field has evolved from manual desk research to
automated signal collection from digital sources: websites, social media,
job postings, financial filings, patents, and news feeds. Modern CI platforms
such as Crayon and Klue~\cite{crayon2024} aggregate these signals into
structured feeds, but they typically stop short of linking external events
to their potential consequences for internal operations, and they do not
integrate with enterprise data systems.

\subsection{NLP for News Understanding and Event Extraction}

Automatically extracting structured information from unstructured news text
is a core NLP task with direct applications in market monitoring. Named
Entity Recognition (NER), Relation Extraction (RE), and Event Detection
pipelines can identify actors, actions, dates, and affected entities in
financial and business news~\cite{devlin2019bert}. LLM-based approaches
have substantially simplified this pipeline: rather than training
specialized extractors, a prompted LLM can parse a news article and output
a structured record describing the event type, entities involved, temporal
context, and assessed business significance~\cite{wei2022chainofthought}.
This approach is used in the present system to convert raw news feeds into
structured entries in the temporal knowledge graph.

\subsection{Job Market Signals as Competitive Indicators}

Competitor job postings have emerged as a reliable leading indicator of
strategic intent~\cite{burning2023jobsignals}. A company hiring aggressively
in a new geographic region, technology domain, or business function signals
expansion plans that may not yet be reflected in public announcements. By
monitoring and analyzing job posting patterns, competitive intelligence
systems can detect strategic pivots, capability investments, and potential
competitive threats weeks or months before they materialize in the market.
The Competitive Intelligence Agent in the present system incorporates job
signal monitoring as one of its primary data streams alongside news and
press releases.

\subsection{Linking Market Events to Internal Operations}

A key differentiator of the present system — and a gap in the existing
literature — is the propagation of external signals through the enterprise
knowledge graph to assess their downstream operational impact. While prior
work has studied market event extraction~\cite{ding2014event} and
portfolio-level financial risk propagation~\cite{acemoglu2012cascade}, the
problem of estimating how a specific external event (a competitor winning
a major contract, an economic downturn in a target sector) affects an
individual organization's projects, workforce allocation, and client
relationships has received limited systematic attention. The present system
addresses this gap by combining GNN-based propagation on the temporal
enterprise knowledge graph with LLM-based causal reasoning and
natural-language recommendation generation.

%─────────────────────────────────────────────────────────────────────────────
\section{LLMs and Knowledge Graphs: Complementary Paradigms}
%─────────────────────────────────────────────────────────────────────────────

Pan et al.~\cite{pan2023llmkg} provided a systematic analysis of the
complementarity between LLMs and knowledge graphs: LLMs bring flexible
natural-language understanding and generative synthesis, while KGs provide
structured, verifiable, and continuously updatable factual grounding.
Three integration modes were identified: (1)~\textit{KG-enhanced LLMs},
where graph-retrieved facts are injected into the LLM context to ground
responses; (2)~\textit{LLM-enhanced KGs}, where language models assist in
KG construction, completion, and enrichment from unstructured sources;
and (3)~\textit{synergistic pipelines} that iterate between both. The
present system employs all three modes: the enterprise KG grounds agent
responses (mode~1), LLM-based extraction from news and documents
continuously enriches the KG (mode~2), and the market analysis pipeline
iterates between LLM analysis and graph updates to progressively refine
the world model (mode~3).

%─────────────────────────────────────────────────────────────────────────────
\section{Project Positioning}
%─────────────────────────────────────────────────────────────────────────────

Table~\ref{tab:sota_positioning} positions the present project relative to
the main research threads reviewed in this chapter.

\begin{table}[H]
\centering
\caption{Positioning of the proposed system relative to the state of the art}
\label{tab:sota_positioning}
\renewcommand{\arraystretch}{1.4}
\resizebox{\textwidth}{!}{%
\begin{tabular}{|p{4.5cm}|p{5cm}|p{6cm}|}
\hline
\textbf{Research area}
  & \textbf{Representative work}
  & \textbf{Use in the proposed system} \\
\hline
LLM reasoning \& tool use
  & ReAct \cite{yao2023react}, CoT \cite{wei2022chainofthought}
  & Domain agents reason and call SQL/graph/API tools in a ReAct loop \\
\hline
LLM task automation
  & HuggingGPT \cite{shen2023hugginggpt}
  & Agents automate leave management, email drafting, and report generation \\
\hline
Multi-agent orchestration
  & AutoGen \cite{wu2023autogen}, LangGraph \cite{langgraph2024}
  & LangGraph orchestrator routes queries to HR/CRM/ERP/Market/CI agents \\
\hline
Retrieval-Augmented Generation
  & Lewis et al. \cite{lewis2020rag}, Self-RAG \cite{asai2023selfrag}
  & Hybrid RAG over internal policies and documents via ChromaDB \\
\hline
Knowledge graphs
  & Hogan et al. \cite{hogan2021knowledge}, GraphRAG \cite{edge2024graphrag}
  & Neo4j enterprise KG linking HR, CRM, ERP, and external market entities \\
\hline
Temporal knowledge graphs
  & Trivedi et al. \cite{trivedi2019dyrep}
  & Temporally annotated KG tracking market event sequences \\
\hline
World model visualization
  & Noy et al. \cite{noy2019industry}
  & Interactive graph canvas mapping internal and external entities \\
\hline
Graph Neural Networks
  & GCN \cite{kipf2017gcn}, GAT \cite{velickovic2018gat}, GraphSAGE \cite{hamilton2017graphsage}
  & GNN-based propagation to assess market event impact on internal operations \\
\hline
LLM + KG integration
  & Pan et al. \cite{pan2023llmkg}
  & Bidirectional enrichment: KG grounds agents; LLMs enrich KG from news \\
\hline
Competitive intelligence
  & Gilad \& Gilad \cite{gilad1989ci}, Crayon \cite{crayon2024}
  & CI Agent monitors news, job postings, and competitor signals \\
\hline
Job signal monitoring
  & Burning Glass \cite{burning2023jobsignals}
  & Competitor hiring trends used as leading strategic indicators \\
\hline
Enterprise data integration
  & Laudon \& Laudon \cite{laudon2020mis}, Text-to-SQL \cite{yu2018spider}
  & Natural language access to ERP, CRM, and HR relational databases \\
\hline
Role-based dashboards
  & Watson \cite{watson2017bi}
  & Role-adaptive dashboards for employee, manager, and executive direction \\
\hline
\end{tabular}%
}
\end{table}

The proposed system distinguishes itself from prior art on four axes.
First, it is the only platform to unify HR, CRM, and ERP domains through
a single multi-agent architecture accessible via natural language, while
simultaneously grounding responses in internal policy documents through RAG.
Second, it links external market intelligence directly to internal operational
data through temporal knowledge graph construction and GNN-based signal
propagation, enabling a form of strategic situational awareness that neither
enterprise AI assistants nor dedicated competitive intelligence platforms
currently provide. Third, its world model visualization bridges the gap
between raw data and strategic comprehension by rendering the enterprise's
full entity-relationship network as an interactive, navigable graph.
Fourth, its architecture is vendor-agnostic and open, built on interoperable
components (FastAPI, LangGraph, Neo4j, ChromaDB, PostgreSQL, React) rather
than on a single vendor's closed ecosystem.

%─────────────────────────────────────────────────────────────────────────────
\section{Conclusion}
%─────────────────────────────────────────────────────────────────────────────

This chapter reviewed the theoretical and technological foundations underpinning
the Intelligent Enterprise Strategic Assistant. Large language models endowed
with tool use and chain-of-thought reasoning provide the generative backbone;
LangGraph's stateful orchestration enables reliable multi-agent coordination
across enterprise domains; hybrid RAG grounds responses in internal
organizational knowledge; and the temporal Neo4j enterprise knowledge graph,
enriched continuously by LLM-based extraction and GNN-based propagation,
connects internal operational entities to external market dynamics. Together,
these building blocks define a coherent technical foundation that no existing
system combines in full. The next chapter presents the system requirements,
the chosen architecture, and the justification of each technology decision
at every layer of the platform.
